\section{Discussion}

\label{sec:discussion}

\subsection{Implications for Conservation}

ZooFed demonstrates that federated learning can bridge the institutional fragmentation of wildlife monitoring data without requiring data centralization.
The 18.7\% average accuracy improvement over local-only models translates to earlier detection of population declines, more accurate occupancy estimates, and better-informed conservation decisions.

\subsection{Limitations}

Several limitations constrain the current system.
First, the HSA architecture requires at least some sites with overlapping sensor types for cross-modal alignment; purely disjoint sensor deployments cannot leverage cross-modal transfer.
Second, the CCDP mechanism requires IUCN Red List assessments, which may lag actual conservation status.
Third, asynchronous participation introduces staleness in gradient updates from intermittently connected sites, partially addressed by learning rate decay for stale updates.

\subsection{Scalability}

The current deployment supports 18 sites.
The NIRA clustering approach scales as $O(K^2)$ for pairwise similarity computation, which becomes expensive beyond $K \approx 100$.
We are investigating approximate nearest-neighbor methods for scalable clustering.

% ============================================================
% 11. Conclusion
% ============================================================
